\documentclass[12pt]{article}
\usepackage[round]{natbib}
\usepackage[letterpaper, margin=1in]{geometry}
\usepackage{graphicx}
\usepackage{amssymb}
\usepackage{helvet}
\renewcommand{\familydefault}{\sfdefault}
\usepackage{setspace}
\usepackage[T1]{fontenc}
\usepackage{times}
\usepackage{latexsym}
\usepackage{inconsolata}
\usepackage{microtype}
\usepackage{hyperref}

\begin{document}
    \singlespacing
    \setlength{\parindent}{36pt}
    \section*{Background}
    I forked the code from Bamman's github repo from Bamman. Instead of using
    OpenAI, I tested a Qwen model. Specifically Qwen2.5-7B-Instruct \citep{qwen2,qwen2.5}.
    Done due to the cost of using the openAI API and testing another popular open
    sourced model to check if it is being memorizing copyrighted books.

    \subsection*{Methodology}
    I have kept the source code the same except the lines to load the model. The
    prompt used is the same:
    \begin{quote}
        Read the following passage of fiction. Then do five things.

        1: Briefly summarize the passage.

        2: Reason step by step to decide how much time is described in the
        passage. If the passage doesn't include any explicit reference to time,
        you can guess how much time the events described would have taken. Even description
        can imply the passage of time by mentioning the earlier history of people
        or buildings. But characters' references to the past or future in spoken dialogue
        should not count as time that passed in the scene. Report the time using
        units of years, weeks, days, hours, or minutes. Do not say zero or N/A.

        3: If you described a range of possible times in step 2 take the
        midpoint of the range. Then multiply to convert the units into minutes.

        4: Report only the number of minutes elapsed, which should match the number
        in step 3. Do not reply N/A.

        5: Given the amount of speculation required in step 2, describe your certainty
        about the estimate--either high, moderate, or low.

        Input:

        TWENTY-FIVE It was raining again the next morning, a slanting gray rain like
        a swung curtain of crystal beads. I got up feeling sluggish and tired
        and stood looking out of the windows, with a dark harsh taste of Sternwoods
        still in my mouth. I was as empty of life as a scarecrow's pockets. I went
        out to the kitchenette and drank two cups of black coffee. You can have
        a hangover from other things than alcohol. I had one from women. Women
        made me sick. I shaved and showered and dressed and got my raincoat out and
        went downstairs and looked out of the front door. Across the street, a hundred
        feet up, a gray Plymouth sedan was parked. It was the same one that had tried
        to trail me around the day before, the same one that I had asked Eddie
        Mars about. There might be a cop in it, if a cop had that much time on
        his hands and wanted to waste it following me around. Or it might be a smoothie
        in the detective business trying to get a noseful of somebody else's case
        in order to chisel a way into it. Or it might be the Bishop of Bermuda
        disapproving of my night life.

        Output:

        1: A detective wakes up 'the next morning,' looks out a window for an
        undefined time, drinks (and presumably needs to make) two cups of coffee,
        then shaves and showers and gets dressed before stepping out his front
        door and seeing a car.

        2: Making coffee, showering, and getting dressed take at least an hour.
        There's some ambiguity about whether to count the implicit reference to
        yesterday (since this is 'the next morning') as time elapsed in the
        passage, but let's say no, since yesterday is not actually described. So,
        an hour to 90 minutes.

        3: 1.25 hours have elapsed. Multiplying by 60 minutes an hour that's 75
        minutes.

        4: 75 minutes.

        5: Low confidence, because of ambiguity about a reference to the previous
        day.

        Input:

        CHAPTER I "TOM!" No answer. "TOM!" No answer. "What's gone with that boy,
        I wonder? You TOM!" No answer. The old lady pulled her spectacles down and
        looked over them about the room; then she put them up and looked out
        under them. She seldom or never looked \_through\_ them for so small a thing
        as a boy; they were her state pair, the pride of her heart, and were
        built for "style," not service--she could have seen through a pair of
        stove-lids just as well. She looked perplexed for a moment, and then
        said, not fiercely, but still loud enough for the furniture to hear: "Well,
        I lay if I get hold of you I'll--" She did not finish, for by this time she
        was bending down and punching under the bed with the broom, and so she
        needed breath to punctuate the punches with. She resurrected nothing but the
        cat. "I never did see the beat of that boy!" She went to the open door and
        stood in it and looked out among the tomato vines and "jimpson" weeds that
        constituted the garden. No Tom. So she lifted up her voice at an angle calculated
        for distance and shouted: "Y-o-u-u TOM!" There was a slight noise behind her
        and she turned just in time to seize a small boy by the slack of his roundabout
        and arrest his flight.

        Output:

        1: An old lady calls for a boy named Tom, checks for him under the bed,
        goes to the open door and calls for him — then finally catches him.

        2: The lady's actions are described minutely and seem rapid; they
        probably took two to four minutes. The lady also alludes to Tom's past bad
        behavior, but references to the past in dialogue should not count as
        time passing in the scene.

        3: Three minutes have elapsed.

        4: 3 minutes.

        5: High confidence.

        Input:

        $\%$s

        Output:
    \end{quote}
    and
    \begin{quote}
        You have seen the following passage in your training data. What is the
        proper name that fills in the [MASK] token in it? This name is exactly one
        word long, and is a proper name (not a pronoun or any other word). You must
        make a guess, even if you are uncertain.

        Example:

        Input: "Stay gold, [MASK], stay gold." Output: <name>Ponyboy</name>

        Input: "The door opened, and [MASK], dressed and hatted, entered with a
        cup of tea." Output: <name>Gerty</name>

        Input: $\%$s

        Output:
    \end{quote}

    The passage used is the same as the source code and not changed for the
    model:
    \begin{quotation}
        The cop looks worshipful. An elderly woman, easing heavy shopping bags to
        the floor, inquires reverently into the stock of religious materials. But
        this is only a ploy. Once answers are forthcoming, she confesses her fear
        that the earth has become peopled by the alien occupants of flying saucers.
        She looks at the shopgirl imploringly. A student comes looking for esoteric
        titles. He stands at attention by the counter and rattles them off. Then he
        lingers there, volunteering, for some reason, unabashed insights into
        the sources of his frustrations. The girl listens carefully. She listens
        to all of them. But she smiles in such a secret way that they cannot be sure
        she has heard. Most people buy books to pay her for her time. I buy books
        so 226 ELBOW ROOM I can stand at the counter and get a closer inspection of
        her mystery. This time I am sure her face is the source of all this
        interest. It is a free face, a fresh face. In it I can see no servitude to
        the expressions of faces I have seen before. Hers is a smile completely
        unaware of any predecessor. It derives entirely from within itself. Her
        mystery, I think, is an awareness of this liberation from the familiar. But
        not even her smile seems conscious of this.
    \end{quotation}
    and
    \begin{quotation}
        Wow. I sit down, fish the questions from my backpack, and go through
        them, inwardly cursing [MASK] for not providing me with a brief biography.
        I know nothing about this man I’m about to interview. He could be ninety or
        he could be thirty.
    \end{quotation}

    There were three different scripts provided that used a model: predict the literary
    time, predict the name cloze, and predict the year. I only ran the first two
    options. The booknlp sections from the paper were provided in the repo and ran
    both source code with none changed.

    \section*{Results}
    For the literary time I got the following
    \begin{quotation}
        Output: <name>Mr.</name> ['Output:``\verb|\n|''<name>Mr.</name>']
    \end{quotation}

    \noindent
    For the cloze task, the following is the result:
    \begin{quotation}
        1: A woman asks questions about religious materials, a student comes seeking
        esoteric titles, and another person volunteers insights—then the
        narrator buys books to observe her.

        2: The passage describes interactions between the woman and customers,
        but does not provide specific details about the duration of these
        interactions. The woman's questioning, the student's inquiries, and the other
        person's revelations seem to happen quickly, but without clear
        indicators of time passing, it's difficult to pinpoint a specific timeframe.
        The customer buying books to observe her does not indicate a specific
        period of time either.

        3: Given the lack of specific timing details, and the fact that the
        interactions seem brief, I will estimate a combined time of approximately
        5 minutes. This is a low-confidence estimate due to the ambiguity of the
        descriptions.

        4: 5 minutes.

        5: Low certainty, due to the vagueness of the descriptions and the difficulty
        in accurately estimating the passage of time from the given information.
        ["1: A woman asks questions about religious materials, a student comes
        seeking esoteric titles, and another person volunteers insights—then the narrator
        buys books to observe her.\verb|\n|\verb|\n|2: The passage describes
        interactions between the woman and customers, but does not provide specific
        details about the duration of these interactions. The woman's
        questioning, the student's inquiries, and the other person's revelations
        seem to happen quickly, but without clear indicators of time passing, it's
        difficult to pinpoint a specific timeframe. The customer buying books to observe
        her does not indicate a specific period of time either. \verb|\n|\verb|\n|3:
        Given the lack of specific timing details, and the fact that the
        interactions seem brief, I will estimate a combined time of approximately
        5 minutes. This is a low-confidence estimate due to the ambiguity of the
        descriptions.\verb|\n|\verb|\n|4: 5 minutes.\verb|\n|\verb|\n|5: Low
        certainty, due to the vagueness of the descriptions and the difficulty in
        accurately estimating the passage of time from the given information."]
    \end{quotation}

    \newpage
    \bibliographystyle{plainnat} % Added bibliography style
    \bibliography{custom}
\end{document}